\documentclass[12pt,a4paper]{article}

\input{/home/tabaka/Documents/GitHub/Defs/defs.tex}

%\pagestyle{empty} 			%usuwa nr strony

\begin{document}
	\begin{mdframed}[style=wzor]
	\large I reguła - dodwanie i odejmowanie ułamków
	\\\\\normalsize 
	Ułamki sprowadzamy do \textbf{wspólnego mianownika}!
	
	Przykłady:
	
	$2\frac{1}{3} + 3 \frac{2}{5} = 2\frac{1\cdot5}{3\cdot5} + 3 \frac{2\cdot3}{5\cdot3} = 2\frac{5}{15} + 3 \frac{6}{15} = 5\frac{11}{15}$\\

	$5\frac{3}{8} - 2 \frac{1}{2} = 5\frac{3}{8} - 2 \frac{1\cdot4}{2\cdot4} = 5\frac{3}{8} - 2 \frac{4}{8} = 3\frac{3}{8}-\frac{4}{8} = 2\frac{3+8}{8} - \frac{4}{8} = 2\frac{11}{8} - \frac{4}{8} = 2\frac{7}{8}$
	
	\end{mdframed}
	
	\begin{mdframed}[style=wzor]
	\large II reguła - mnożenie i dzielenie ułamków
	\\\\\normalsize
	Ułamki sprowadzamy do \textbf{ułamków niewłaściwych}!
	
	Przykłady:
	
	$5\cdot 2\frac{1}{5} = \frac{5}{1}\cdot\frac{11}{5}= \frac{{}^1 \cancel{5}}{1}\cdot\frac{11}{\cancel{5}_1}= \frac{11}{1} = 11$\\
	
	$1\frac{2}{3}:\frac{10}{21} = \frac{5}{3}\cdot \frac{21}{10} = \frac{{}^1\cancel{5}}{{}_1\cancel{3}}\cdot \frac{\cancel{21}^7}{\cancel{10}_2} = \frac{7}{2} = 3\frac{1}{2}$
	
	\end{mdframed}
	
	\vspace{1.5cm}
	
	\begin{mdframed}[style=wzor]
		\large Jeśli mamy dodawanie lub udejmowanie \textbf{NIE SKRACAMY} minusów.
		\\\\\normalsize
		Przykłady:
		
		$4 + 5 = 9$\\
		
		$9 + (-5) = 9 - 5 = 4$\\
		
		$-9 + (-5) = -14 $\\
		
		$-9 + 5 = -4$
		
	\end{mdframed}
	
	\begin{mdframed}[style=wzor]
		\large Jeśli mamy mnożenie lub dzielenie to \textbf{SKRACAMY} minusy.
		\\\\\normalsize
		Przykłady:
		
		$-5\cdot (-4) = +5 \cdot (+4) = 20$\\
		
		$3 \cdot (-7) = 21$\\
		
		$-6 : (-2) = 6 : 2 = 3$\\
		
		$(-2) \cdot (-1) \cdot (-5) \cdot (-4) = (+2) \cdot (+1) \cdot (-5) \cdot (-4)= (+2) \cdot (+1) \cdot (+5) \cdot (+4) = 2\cdot1\cdot5\cdot4 = 40$
	\end{mdframed}

\end{document}